\chapter{Low Latency}\label{ch:iv3:low-latency}

\begin{ivtrains}
Reading code the way the machine and the language rules read it: what a listing really does, why one version is slower, and which guarantee a piece
of concurrent code is relying on. Most questions name the compiler and processor they depend on, because the honest answer often changes with them. The
later ones ask you to write the structure and then defend every line.
\end{ivtrains}

\begin{question}[$\star$ \iqroles{QD, Dev, LLE} \iqfirm{any} \iqlevel{junior}]\label{iq:iv3:low-latency:1}
Compiled with g++ for x86-64 Linux, what does \cref{lst:iv3:ll:layout} print? Which part of the answer does the C++ standard fix, and which part belongs to the platform?
\omcode{code/interviews-3/04-low-latency/cpp/layout.cpp}{4}{12}{Three layouts.}\label{lst:iv3:ll:layout}
\end{question}

\begin{question}[$\star$ \iqroles{QD, Dev, LLE} \iqfirm{any} \iqlevel{junior}]\label{iq:iv3:low-latency:2}
\Cref{lst:iv3:ll:spinflag} is compiled with g++ 11.4 on x86-64, \texttt{-std=c++20 -O2}, and its author expects 42. Is the program defined? Can you say what it prints? Does declaring the flag \texttt{volatile bool} fix it, and if not, what does?
\omcode{code/interviews-3/04-low-latency/cpp/spin_flag.cpp}{5}{18}{A flag and a value.}\label{lst:iv3:ll:spinflag}
\end{question}

\begin{question}[$\star$ \iqroles{QD, Dev, LLE} \iqfirm{crypto, market maker} \iqlevel{junior}]\label{iq:iv3:low-latency:3}
With rustc 1.97 on x86-64, what does \cref{lst:iv3:ll:sizes} print, line by line? Which of the sizes of the \texttt{Option} types does Rust promise, and which merely happen to come out that way?
\omcode{code/interviews-3/04-low-latency/rust/src/bin/sizes.rs}{1}{15}{Sizes of ten types.}\label{lst:iv3:ll:sizes}
\end{question}

\begin{question}[$\star$ \iqroles{QD, Dev, LLE} \iqfirm{crypto, market maker} \iqlevel{junior}]\label{iq:iv3:low-latency:4}
\Cref{lst:iv3:ll:spawn} compiles and prints 300. Delete the word \texttt{move} on line 5. What happens, and at which stage? The thread is joined on the very next line, so why is that not enough? Give a fix that needs no reference count and leaves \texttt{prices} usable in \texttt{main} afterwards.
\omcode{code/interviews-3/04-low-latency/rust/src/bin/spawn_borrow.rs}{1}{7}{A worker thread.}\label{lst:iv3:ll:spawn}
\end{question}

\begin{question}[$\star$ \iqroles{LLE, Dev, QD, QT} \iqfirm{market maker} \iqlevel{junior}]\label{iq:iv3:low-latency:5}
Take these round numbers: a 4 GHz core, 2 instructions per cycle, and a cache miss all the way to main memory costing 90 ns. You have a budget of 1 microsecond from packet to order. How many cycles and instructions is that? How many misses to main memory fill the whole budget? What share of the budget goes on 5 misses, and how many instructions could have run in that time?
\end{question}

\begin{question}[$\star\star$ \iqroles{LLE, Dev, QD} \iqfirm{market maker, hedge fund} \iqlevel{junior, experienced}]\label{iq:iv3:low-latency:6}
In \cref{lst:iv3:ll:counters}, eight threads each run \texttt{worker} with their own \texttt{t}, on a machine with 64-byte cache lines. No two threads touch the same counter and every increment is atomic, yet it scales badly: the more threads, the lower the total rate. Explain the mechanism, fix it, and say what the fix costs. How would you confirm the diagnosis?
\omcode{code/interviews-3/04-low-latency/cpp/counters.cpp}{9}{17}{One counter per thread.}\label{lst:iv3:ll:counters}
\end{question}

\begin{question}[$\star\star$ \iqroles{LLE, Dev} \iqfirm{market maker, hedge fund} \iqlevel{junior, experienced}]\label{iq:iv3:low-latency:7}
Compile \cref{lst:iv3:ll:atomicops} with g++ 11.4, \texttt{-std=c++20 -O2}, for x86-64. Which of these eight functions contain only an ordinary move or add, and which need an instruction that is locked or otherwise more expensive? Explain the pattern. What changes on a weakly ordered processor?
\omcode{code/interviews-3/04-low-latency/cpp/atomic_ops.cpp}{3}{13}{Eight small functions.}\label{lst:iv3:ll:atomicops}
\end{question}

\begin{question}[$\star\star$ \iqroles{LLE, Dev, QD} \iqfirm{any} \iqlevel{junior, experienced}]\label{iq:iv3:low-latency:8}
In \cref{lst:iv3:ll:lists}, \texttt{values} starts as $0, 1, \dots, n-1$, with $n$ two million, and the two lists end up holding the same values in the same order. The same function, which walks a list and adds up its values, takes many times longer on \texttt{b} than on \texttt{a}. Why? How would you confirm your explanation without any hardware counters, and what would you change?
\omcode{code/interviews-3/04-low-latency/cpp/list_order.cpp}{41}{45}{Two equal lists.}\label{lst:iv3:ll:lists}
\end{question}

\begin{question}[$\star\star$ \iqroles{QD, Dev, LLE} \iqfirm{any} \iqlevel{junior, experienced}]\label{iq:iv3:low-latency:9}
\Cref{lst:iv3:ll:countnew} replaces global \texttt{operator new} with one that counts its calls. With g++ 11.4 and its standard library, \texttt{-std=c++20 -O2}, x86-64, what count does each line print? Which of these numbers could a different standard library change?
\omcode{code/interviews-3/04-low-latency/cpp/count_new.cpp}{10}{47}{Counting allocations.}\label{lst:iv3:ll:countnew}
\end{question}

\begin{question}[$\star\star$ \iqroles{LLE, Dev, QD} \iqfirm{market maker, crypto} \iqlevel{junior, experienced}]\label{iq:iv3:low-latency:10}
A packet sits in a byte buffer, with a header starting 3 bytes in. With g++ 11.4 on x86-64, \texttt{-std=c++20 -O2}, \cref{lst:iv3:ll:parse} prints the right value twice. Are both functions correct under C++20? What do you expect the machine code of each to look like?
\omcode{code/interviews-3/04-low-latency/cpp/parse_header.cpp}{6}{27}{Two ways to read a timestamp.}\label{lst:iv3:ll:parse}
\end{question}

\begin{question}[$\star\star$ \iqroles{LLE, Dev, QD, QR} \iqfirm{any} \iqlevel{junior, experienced}]\label{iq:iv3:low-latency:11}
\Cref{lst:iv3:ll:benchbad}, built with g++ 11.4, \texttt{-O2}, on x86-64, prints about 70 nanoseconds for a billion iterations. Why? Fix it, and list what else is wrong with it as a way to measure the latency of \texttt{work}. If you then read the time-stamp counter directly with \texttt{rdtsc} instead of a library clock, what new mistake can you make?
\omcode{code/interviews-3/04-low-latency/cpp/bench_bad.cpp}{5}{16}{A benchmark.}\label{lst:iv3:ll:benchbad}
\end{question}

\begin{question}[$\star\star$ \iqroles{LLE, Dev} \iqfirm{market maker, hedge fund} \iqlevel{junior, experienced}]\label{iq:iv3:low-latency:12}
With g++ 11.4, \texttt{-std=c++20 -O2}, on x86-64, \cref{lst:iv3:ll:publish} prints 42 in every run you try. Is the program correct under the C++20 memory model? Would anything change on ARM or POWER? Fix it with the weakest orders that work.
\omcode{code/interviews-3/04-low-latency/cpp/publish.cpp}{5}{20}{Publishing a value.}\label{lst:iv3:ll:publish}
\end{question}

\begin{question}[$\star\star$ \iqroles{LLE, Dev} \iqfirm{market maker, hedge fund} \iqlevel{junior, experienced}]\label{iq:iv3:low-latency:13}
With g++ 11.4, \texttt{-O2}, on x86-64, four threads each lock this spinlock and increment a plain \texttt{long} counter 20\,000 times, and the total comes out right. \Cref{lst:iv3:ll:spinlock} still has one fault of correctness and one of performance under contention. Find both and fix them.
\omcode{code/interviews-3/04-low-latency/cpp/spinlock.cpp}{6}{15}{A spinlock.}\label{lst:iv3:ll:spinlock}
\end{question}

\begin{question}[$\star\star$ \iqroles{LLE, Dev, QD} \iqfirm{market maker, crypto} \iqlevel{experienced}]\label{iq:iv3:low-latency:14}
You want a thread to publish a best bid and offer that others read atomically. \Cref{lst:iv3:ll:quote} is built with g++ 11.4, \texttt{-std=c++20 -O2 -latomic}, on an x86-64 processor that supports the 16-byte compare-and-swap. It reports the 8-byte type as lock-free, and the 16- and 24-byte types as not lock-free. What does that mean for a hot path, and what happens if you leave out \texttt{-latomic}? What would you do instead?
\omcode{code/interviews-3/04-low-latency/cpp/quote_atomic.cpp}{4}{20}{Three sizes of atomic.}\label{lst:iv3:ll:quote}
\end{question}

\begin{question}[$\star\star$ \iqroles{QD, Dev, LLE} \iqfirm{crypto, market maker} \iqlevel{junior, experienced}]\label{iq:iv3:low-latency:15}
In Rust, a type that is \texttt{Send} may be moved to another thread and a type that is \texttt{Sync} may be shared by reference between threads. For each of these types, say whether it is \texttt{Send} and whether it is \texttt{Sync}, with the reason: \texttt{Rc<u64>}, \texttt{Arc<u64>}, \texttt{Arc<Cell<u64>{}>}, \texttt{Cell<u64>}, \texttt{RefCell<u64>}, \texttt{Mutex<Cell<u64>{}>}, \texttt{Mutex<Rc<u64>{}>}, \texttt{\&\textquotesingle{}static Cell<u64>}, \texttt{*mut u64}, \texttt{MutexGuard<\textquotesingle{}static, u64>}, \texttt{AtomicU64}.
\end{question}

\begin{question}[$\star\star$ \iqroles{QD, Dev, LLE} \iqfirm{crypto, market maker} \iqlevel{junior, experienced}]\label{iq:iv3:low-latency:16}
\Cref{lst:iv3:ll:lost} contains no \texttt{unsafe} and compiles with no warnings under rustc 1.97. What can it print? Does it have a data race, or undefined behaviour? What is the smallest final value the rules allow, and how would you fix the program?
\omcode{code/interviews-3/04-low-latency/rust/src/bin/lost_update.rs}{1}{18}{A shared counter, in safe code.\protect\hfil\protect\break}\label{lst:iv3:ll:lost}
\end{question}

\begin{question}[$\star\star$ \iqroles{QD, Dev, LLE} \iqfirm{crypto, market maker} \iqlevel{junior, experienced}]\label{iq:iv3:low-latency:17}
A global allocator counts every call to allocate or to reallocate, and \texttt{show} in \cref{lst:iv3:ll:countalloc} prints how many calls its closure made. With rustc 1.97 in a release build, and given that a \texttt{Vec} growing from empty starts at capacity 4 and doubles, what does each line print? Which of the numbers does the documentation promise?
\omcode{code/interviews-3/04-low-latency/rust/src/bin/count_alloc.rs}{25}{53}{Counting allocator calls.}\label{lst:iv3:ll:countalloc}
\end{question}

\begin{question}[$\star\star$ \iqroles{QD, Dev, LLE} \iqfirm{crypto, market maker} \iqlevel{experienced}]\label{iq:iv3:low-latency:18}
Compile \cref{lst:iv3:ll:bounds} with rustc 1.97, \texttt{-C opt-level=3}, on x86-64. An index that may be out of range normally needs a comparison and a branch to a panic. Which of these seven functions still have a bounds check on the indexing, which have some other check once at the start, and which have none at all? Is any of this guaranteed?
\omcode{code/interviews-3/04-low-latency/rust/src/bounds.rs}{1}{36}{Seven ways to index.}\label{lst:iv3:ll:bounds}
\end{question}

\begin{question}[$\star\star$ \iqroles{LLE, Dev, QD} \iqfirm{market maker, hedge fund} \iqlevel{junior, experienced}]\label{iq:iv3:low-latency:19}
A process reads random entries of an 8 GiB table. Assume the translation cache (the TLB) holds 2048 entries, a miss in it costs a 30 ns page-table walk, and the walk costs the same at any page size. For a page size of 4 KiB, then of 2 MiB, estimate how much of the table the TLB covers, the miss rate, and the average translation cost per read. How much memory do the last-level page-table entries take at each page size, at 8 bytes per entry? What are the catches of huge pages?
\end{question}

\begin{question}[$\star\star\star$ \iqroles{LLE, Dev, QD} \iqfirm{market maker, hedge fund} \iqlevel{experienced}]\label{iq:iv3:low-latency:20}
Two \texttt{std::atomic<int>} variables \texttt{x} and \texttt{y} start at 0. Thread 1 does \texttt{x.store(1); r1 = y.load();}. Thread 2 does \texttt{y.store(1); r2 = x.load();}. Can both registers end at 0 if (a) every operation is relaxed, (b) stores are release and loads are acquire, (c) every operation is sequentially consistent? For each, answer for the C++20 language and for x86-64 hardware. Which instruction makes the difference in (c), and how would you show the result on a real machine?
\end{question}

\begin{question}[$\star\star\star$ \iqroles{LLE, Dev} \iqfirm{market maker, hedge fund, crypto} \iqlevel{junior, experienced}]\label{iq:iv3:low-latency:21}
On a whiteboard, in C++20, write a fixed-capacity ring buffer for exactly one producer thread and one consumer thread, with no locks. \texttt{push} returns false when the buffer is full and \texttt{pop} returns false when it is empty, and the capacity is a power of two. Then defend the memory order of every atomic operation: what can go wrong if you weaken it?
\end{question}

\begin{question}[$\star\star\star$ \iqroles{LLE, Dev} \iqfirm{market maker, hedge fund} \iqlevel{experienced}]\label{iq:iv3:low-latency:22}
\Cref{lst:iv3:ll:adjacent} is a correct single-producer single-consumer queue, with the two threads on different cores. Putting \texttt{alignas(64)} on \texttt{head}, \texttt{tail} and \texttt{slots} makes it noticeably faster, and one further change, with every memory order left as it is, makes it faster again. Explain both mechanisms, name the further change, and say what it costs. How would you confirm each effect?
\omcode{code/interviews-3/04-low-latency/cpp/spsc_layout.cpp}{11}{29}{A queue with its fields side by side.}\label{lst:iv3:ll:adjacent}
\end{question}

\begin{question}[$\star\star\star$ \iqroles{LLE, Dev, QD} \iqfirm{market maker, crypto} \iqlevel{experienced}]\label{iq:iv3:low-latency:23}
\Cref{lst:iv3:ll:wrap} is shown single-threaded to keep it short. Its indices run freely and rely on unsigned wrap-around, and every test passes. At a million messages a second, it can start to lose data after a while. After how long, under exactly what condition, and what is the smallest fix? Is \texttt{tail - head} the problem?
\omcode{code/interviews-3/04-low-latency/cpp/wrap.cpp}{4}{21}{A ring with free-running indices.}\label{lst:iv3:ll:wrap}
\end{question}

\begin{question}[$\star\star\star$ \iqroles{LLE, Dev, QD, QR} \iqfirm{market maker, hedge fund} \iqlevel{experienced}]\label{iq:iv3:low-latency:24}
A hot thread must record each latency, an integer number of nanoseconds from 1 ns to about 1 s, so that any percentile can be read later to within 2\%. It may not allocate, may not use floating point, and must take constant time per record. Design the histogram: how many buckets, how much memory, and what is the indexing function in C++20? Why not one bucket per nanosecond?
\end{question}

\begin{question}[$\star\star\star$ \iqroles{LLE, Dev} \iqfirm{market maker, hedge fund} \iqlevel{experienced}]\label{iq:iv3:low-latency:25}
\Cref{lst:iv3:ll:firsttouch} runs on x86-64 Linux with 4 KiB pages, glibc, and transparent huge pages not forced on. Roughly what five numbers does it print? Use your answer to explain why the first messages after a process starts can be slow, and what you would do at start-up. On a two-socket machine, where does the vector's memory end up if \texttt{main} builds it and a thread pinned to the other socket then uses it?
\omcode{code/interviews-3/04-low-latency/cpp/first_touch.cpp}{16}{35}{Minor page faults, counted.}\label{lst:iv3:ll:firsttouch}
\end{question}

\begin{question}[$\star\star\star$ \iqroles{LLE, Dev, QD, QR, QT} \iqfirm{market maker, hedge fund} \iqlevel{junior, experienced}]\label{iq:iv3:low-latency:26}
A load tool sends one request every 2 ms, but waits for each reply before sending the next, and records the time from send to reply. The server is a single first-in first-out queue that takes 1 ms per request. During a 100-second run it freezes once, for exactly 1 second. Estimate by hand what the tool reports for its 99th and 99.9th percentiles and its maximum. Then estimate the same percentiles for a client that keeps to the 2 ms schedule no matter what, measuring from each request's scheduled send time. Which do you believe, and how do you change the tool?
\end{question}

\begin{question}[$\star\star\star$ \iqroles{LLE, Dev} \iqfirm{market maker, hedge fund} \iqlevel{experienced}]\label{iq:iv3:low-latency:27}
A 10 Gb/s Ethernet port receives market data in a burst at full line rate, in the smallest frames: 64 bytes, plus 20 bytes of preamble and gap on the wire. How long do you have per packet, and how many cycles at 3 GHz? Assume a receive system call costs 500 ns and the kernel network stack another 150 ns per packet. Is one call per packet enough? One call per 32 packets? Now choose among batched sockets with busy polling, an \texttt{AF\_XDP} socket and DPDK, when the same port must also keep carrying ordinary management traffic and the team wants to keep using standard kernel tools. What does each option give up?
\end{question}

\begin{question}[$\star\star\star$ \iqroles{LLE, Dev} \iqfirm{market maker, hedge fund} \iqlevel{experienced}]\label{iq:iv3:low-latency:28}
On Linux, a thread pinned to a core removed from the scheduler with \texttt{isolcpus} spins on a queue and makes no system calls. Most iterations take 50 ns, but every few seconds one takes tens of microseconds, and now and then one takes over a millisecond. What can still interrupt that core? Give the order in which you would check, the tool or file for each check, and the remedy.
\end{question}

\begin{question}[$\star\star\star$ \iqroles{QD, Dev, LLE} \iqfirm{crypto, market maker} \iqlevel{experienced}]\label{iq:iv3:low-latency:29}
With rustc 1.97 on x86-64, what do the first two lines of output from \cref{lst:iv3:ll:header} show, and which of those four sizes are guaranteed? \texttt{ts\_cast} is then called on a slice starting 3 bytes into an 8-aligned buffer. What happens in a debug build and in a release build, and what does the language say? For a value \texttt{p} of type \texttt{Packed}, why does writing \texttt{\{ p.ts \}} compile, when taking \texttt{\&p.ts} does not? How should \texttt{ts\_cast} be written?
\omcode{code/interviews-3/04-low-latency/rust/src/bin/header.rs}{4}{31}{Layouts and a cast.}\label{lst:iv3:ll:header}
\end{question}

\begin{question}[$\star\star\star$ \iqroles{QD, Dev, LLE} \iqfirm{crypto, market maker} \iqlevel{experienced}]\label{iq:iv3:low-latency:30}
\Cref{lst:iv3:ll:mailbox} is a one-item mailbox that its author describes as safe to use from safe code. Find every way that code containing no \texttt{unsafe} at all can use this API to cause undefined behaviour, and any way it can misbehave without undefined behaviour. Fix it, changing the API where you must.
\omcode{code/interviews-3/04-low-latency/rust/src/bin/mailbox.rs}{7}{34}{A one-item mailbox.}\label{lst:iv3:ll:mailbox}
\end{question}

\begin{question}[$\star\star\star\star$ \iqroles{LLE, Dev} \iqfirm{market maker, hedge fund} \iqlevel{experienced}]\label{iq:iv3:low-latency:31}
\Cref{lst:iv3:ll:seqlock} is a reader-writer scheme for a quote: one writer, many readers, and a reader retries if it saw the sequence number odd or changed. Is it correct under the C++20 memory model? Name every way it can fail, whether or not x86-64 hides it. Rewrite both sides, and defend every order you use. What does the scheme cost the writer and the readers, compared with a mutex?
\omcode{code/interviews-3/04-low-latency/cpp/seqlock_naive.cpp}{9}{29}{A sequence number and a quote.}\label{lst:iv3:ll:seqlock}
\end{question}

\begin{question}[$\star\star\star\star$ \iqroles{LLE, Dev} \iqfirm{market maker, hedge fund} \iqlevel{experienced}]\label{iq:iv3:low-latency:32}
Many threads call \texttt{push} and \texttt{pop} on the stack in \cref{lst:iv3:ll:treiber}. It has two separate faults, and each can corrupt memory or crash under g++ on x86-64, whatever the memory orders. Give, for each, the exact order of events between two threads that shows it. Then compare the ways to fix it: what does each one cost?
\omcode{code/interviews-3/04-low-latency/cpp/treiber.cpp}{6}{31}{A stack built on compare-and-swap.}\label{lst:iv3:ll:treiber}
\end{question}

\begin{question}[$\star\star\star\star$ \iqroles{LLE, Dev} \iqfirm{market maker, hedge fund} \iqlevel{experienced}]\label{iq:iv3:low-latency:33}
Call an algorithm lock-free if, whichever threads are suspended, some running thread always completes its operation in a finite number of its own steps. \Cref{lst:iv3:ll:mpmc} shows the push half of a bounded queue for many producers and many consumers. It uses no mutex, and \texttt{try\_pop}, not shown, is its mirror image: it takes a position only when that slot's \texttt{seq} says the item is ready. Suppose one producer is suspended just after its compare-and-swap, starting on line 28, succeeds, on an 8-slot queue. What can the other threads still do, and what do they see? Is the queue lock-free in the sense above? What is the worst-case latency it can impose on a consumer, and how would you redesign if that is unacceptable?
\omcode{code/interviews-3/04-low-latency/cpp/mpmc.cpp}{7}{40}{The push half of a bounded queue.}\label{lst:iv3:ll:mpmc}
\end{question}

\begin{question}[$\star\star\star\star$ \iqroles{QD, Dev, LLE} \iqfirm{crypto, market maker} \iqlevel{experienced}]\label{iq:iv3:low-latency:34}
In Rust, with no external crates, design a bounded single-producer single-consumer queue whose API makes these three mistakes fail to compile: two threads pushing at once, popping again while still holding a reference to a peeked item, and sending a non-thread-safe item type across it. Give the types and signatures, write \texttt{push} and a zero-copy \texttt{peek}, and state which \texttt{unsafe impl} lines you need and why. Which errors will the compiler still not catch for you?
\end{question}

\begin{question}[$\star\star\star\star$ \iqroles{QD, Dev, LLE} \iqfirm{crypto, market maker} \iqlevel{experienced}]\label{iq:iv3:low-latency:35}
\Cref{lst:iv3:ll:myarc} is a hand-written atomically reference-counted pointer in which every count operation is \texttt{Relaxed}. Which operations are correct as they stand, and which are not? Give an order of events that goes wrong, and fix it with the weakest orders that work. Why must both bounds \texttt{T: Send + Sync} be there, and what does your fix cost on x86-64 and on a weakly ordered processor?
\omcode{code/interviews-3/04-low-latency/rust/src/my_arc_bad.rs}{1}{37}{A reference-counted pointer.}\label{lst:iv3:ll:myarc}
\end{question}
